\documentclass{article}
\usepackage{amsmath,graphicx}
\newcommand{\myvec}[1]{\mathbf{#1}}

\title{Mock Rectified Flow}
\author{Xingchao Liu \and Chengyue Gong}
\date{}

\begin{document}
\maketitle

\begin{abstract}
We present rectified flow, a mock abstract that must never leak into body blocks.
\end{abstract}

\section{Introduction}
\label{sec:intro}
Rectified flow transports data between two distributions~\cite{liu2022flow} via straight
paths, see~\eqref{eq:ode}. See Figure~\ref{fig:overview} and Table~\ref{tab:results} for
details. \emph{Straightness} reduces sampling cost. A footnote about causal
dynamics\footnote{This is a footnote about causal dynamics.} follows here.

\begin{equation}
\label{eq:ode}
\mathrm{d}Z_t = v(Z_t, t)\,\mathrm{d}t
\end{equation}

\subsection{Reflow}
\label{sec:reflow}
The reflow operator straightens paths further, described by $y = f(x)$ inline math.

\begin{figure}[t]
\centering
\includegraphics{x1.png}
\caption{Overview of the rectified flow process.}
\label{fig:overview}
\end{figure}

\section{Method}
\label{sec:method}
The method section describes the training procedure (see \texttt{reflow.py}), see
Section~\ref{sec:intro} again.

\begin{align}
\label{eq:group}
v^{*} &= \arg\min_{v}\, \mathbb{E}\,\|(X_1 - X_0) - v\|^2 \\
Z_0 &= X_0
\end{align}

\begin{table}[t]
\centering
\begin{tabular}{ll}
Method & Score \\
Ours & 0.99 \\
\end{tabular}
\caption{Results comparison.}
\label{tab:results}
\end{table}

\begin{itemize}
\item First item referencing Section~\ref{sec:intro}.
\item Second item about reflow.
\end{itemize}

\begin{theorem}
\label{thm:main}
Rectified flow paths are optimal transport maps under mild conditions.
\end{theorem}

\begin{algorithm}
\caption{Rectified Flow Sampling}
\label{alg:sampling}
\begin{algorithmic}
\FOR{$i = 0$ to $N$}
\STATE for i in range(N): update state
\ENDFOR
\end{algorithmic}
\end{algorithm}

\begin{verbatim}
pip install rectified-flow
\end{verbatim}

See our code at \url{https://github.com/gnobitab/RectifiedFlow} for details.

\begin{quote}
Straight paths are the shortest paths between two distributions.
\end{quote}

\appendix
\section{Proofs}
\label{sec:proofs}
Appendix proof text that should not be auto translated at all, referencing~\ref{sec:method}.


\begin{thebibliography}{9}
\bibitem{liu2022flow} Xingchao Liu, Chengyue Gong, Qiang Liu. \emph{Flow Straight and Fast:
Learning to Generate and Transfer Data with Rectified Flow}. arXiv:2209.03003, 2022.
\bibitem{song2020ddpm} Jiaming Song, Chenlin Meng, Stefano Ermon. Denoising Diffusion
Implicit Models. \url{https://doi.org/10.48550/arXiv.2010.02502}, 2020.
\end{thebibliography}

\end{document}
